%==============================================================================
%  Overleaf-ready manuscript (v2, repositioned after hard reject)
%
%  How to use:
%   - This compiles standalone on Overleaf with pdfLaTeX.
%   - Red text  \tbd{...}   = a value/table you must fill in AFTER re-running the
%                             closed-loop calibration on the TGSIM data.
%   - Blue text \prelim{...}= a number already computed by the analysis code but
%                             from the OLD (pre-fix) class means; regenerate it
%                             after recalibration via analysis/run_emergent_analysis.py
%   - To drop into the official TRB/TRR template, copy the sections below.
%==============================================================================
\documentclass[11pt]{article}

\usepackage[letterpaper,margin=1in]{geometry}
\usepackage{mathptmx}            % Times-like font (TRR style)
\usepackage{amsmath,amssymb}
\usepackage{graphicx}
\usepackage{booktabs}
\usepackage{multirow}
\usepackage{array}
\usepackage{setspace}
\usepackage[dvipsnames]{xcolor}
\usepackage[hidelinks]{hyperref}
\usepackage{lineno}
\usepackage{caption}

% --- helper macros for drafting -------------------------------------------------
\newcommand{\tbd}[1]{\textcolor{red}{[#1]}}        % fill after re-run
\newcommand{\prelim}[1]{\textcolor{blue}{#1}}      % preliminary (old means)
\newcommand{\dv}{\Delta v}

\title{\textbf{Beyond Calibration: Regime-Dependent Car-Following of
Autonomous versus Human-Driven Vehicles and Its Emergent String-Stability
and Safety Consequences}}

\author{Pedram Beigi, Mohammad Emad Rashidi, Nachuan Li, Shayan Bafandkar,\\
Dana Monzer, John Hourdos, Hani Mahmassani, Alireza Talebpour, Samer H. Hamdar}
\date{}

\begin{document}
\linenumbers
\maketitle

%================================================================= ABSTRACT
\begin{abstract}
\noindent
This study moves beyond descriptive parameter comparison to examine how the
longitudinal behavior of small vehicles, large vehicles, and autonomous vehicles
(AVs) translates into emergent traffic-flow and safety outcomes. Using
high-resolution TGSIM trajectory data, we calibrate four car-following models
spanning the principal modeling paradigms---the physics-based Intelligent Driver
Model (IDM), the behaviorally motivated Prospect-Theory (PT) model, the
optimal-velocity relative-velocity (OVRV) model, and the safety-based Gipps
model, complemented by an ACC-oriented improved IDM as an AV-control
surrogate---at the individual leader--follower level using a \emph{closed-loop}
trajectory simulation, in which the spacing that drives the model is recomputed
from the simulated follower position at every step. We make three contributions. First, we stratify
car-following by traffic regime (free-flow, synchronized, and stop-and-go) and
show that class differences which are diluted in the pooled analysis emerge
within specific regimes. Second, we assess \emph{out-of-sample} predictive
accuracy and \emph{practical identifiability}: using the dispersion of
multi-start calibration and one-at-a-time profiles, we show that IDM parameters
related to spacing and desired speed are well identified and transferable,
whereas several PT parameters are practically non-identifiable from short
trajectories---reframing the limited PT differentiation as an identifiability
property rather than an absence of behavioral difference. Third, we map the
calibrated class- and regime-specific parameters to emergent dynamics through an
analytical linear string-stability criterion and mixed-autonomy platoon
simulations, quantifying string amplification, surrogate safety measures
(time-to-collision exposure and required deceleration), and throughput across AV
market-penetration rates. The results show that the calibrated AV behavior trades
modest throughput gains against reduced string stability and shrinking safety
margins at high penetration, with consequences that are strongly regime
dependent. The framework yields class- and regime-specific parameters that are
directly usable as priors for mixed-autonomy simulation and for lane-management
and AV-integration policy.

\vspace{0.5em}
\noindent\textbf{Keywords:} Car-following, Autonomous Vehicles, Prospect Theory,
Traffic Regimes, String Stability, Parameter Identifiability, Naturalistic
Trajectory Data.
\end{abstract}

%================================================================= INTRODUCTION
\section{Introduction}
Microscopic simulation is a pivotal tool for the analysis and design of traffic
systems, allowing controlled study of new technologies without disrupting real
traffic~\cite{chung2009}. Car-following (CF) models are the core of microscopic
simulation, reproducing the longitudinal motion of a vehicle behind its
leader~\cite{zheng2023}. Despite decades of development, our understanding of CF
behavior remains incomplete, constrained by data, by the difficulty of model
calibration and validation, and---most acutely---by the heterogeneity of driving
agents, which now includes autonomous vehicles (AVs)~\cite{zheng2016}.

AVs are expected to increase throughput and improve stability, but the empirical
evidence is mixed: commercial adaptive-cruise-control systems have been shown to
be string \emph{unstable} in several field experiments~\cite{gunter2020}. While
AV longitudinal control has been studied extensively, there remains limited
empirical clarity on (i) how AV behavior differs from human-driven vehicles
(HDVs) \emph{across traffic regimes}, (ii) whether the calibrated parameters are
\emph{identifiable and transferable}, and (iii) what these parameters imply for
\emph{emergent} flow and safety in mixed traffic. The first version of this study
calibrated IDM and PT per leader--follower pair and compared the resulting
parameters statistically. That analysis, while internally consistent, shared the
limitations of much of the literature: it pooled all traffic conditions, reported
only in-sample fit, and stopped at the parameters themselves without connecting
them to emergent behavior.

This paper addresses those gaps directly. We use the naturalistic TGSIM
dataset~\cite{ammourah2024} and classify the observed followers into
(i)~\emph{small} vehicles (passenger cars, SUVs, pickups), (ii)~\emph{large}
vehicles (trucks and buses), and (iii)~\emph{autonomous} vehicles. Our
contributions are:
\begin{enumerate}
\item \textbf{Regime-stratified calibration and comparison.} Each car-following
episode is labeled by traffic regime (free-flow, synchronized, stop-and-go), and
class differences are tested both pooled and within regime, revealing
regime-specific behavioral distinctions that pooling conceals.
\item \textbf{Validation and practical identifiability.} We evaluate
out-of-sample prediction (calibrate on part of an episode, predict the rest) and
quantify practical identifiability from multi-start dispersion and profile
analysis, establishing which parameters are trustworthy and explaining the
limited PT differentiation as non-identifiability.
\item \textbf{Emergent string-stability and safety analysis.} We translate the
calibrated class/regime parameters into analytical linear string stability and
mixed-autonomy platoon simulations, quantifying string amplification, surrogate
safety, and throughput as functions of AV market penetration.
\end{enumerate}

%================================================================= LITERATURE
\section{Literature Review}
\subsection{Calibration of car-following models}
Among physics-based models, the IDM~\cite{treiber2000} reproduces realistic
fundamental diagrams and stable/unstable regions and consistently shows low
calibration error relative to alternatives~\cite{zhu2018,pourabdollah2018}.
Trajectory-based calibration is now standard~\cite{punzo2005}, but a recurrent
methodological pitfall is the choice of simulation scheme. In an \emph{open-loop}
(one-step) scheme the model is fed observed states at each step; in a
\emph{closed-loop} scheme the spacing and speed are integrated forward and the
spacing that drives the model is recomputed from the simulated position. Punzo and
co-authors emphasize that closed-loop, spacing-based objectives are required for
trajectory calibration to be sensitive to the parameters and to avoid optimistic,
non-identifiable fits~\cite{punzo2012,montanino2015}. We adopt a closed-loop
scheme for \emph{both} models in this study.

\subsection{Behavioral models and Prospect Theory}
Psychologically grounded models capture driver decision-making under
uncertainty. The Prospect-Theory CF model~\cite{hamdar2008rec,hamdar2015}
represents acceleration as a risk-taking choice combining subjective utilities
for accelerations and disutilities for crash risk, calibrated with a genetic
algorithm to account for inter-driver heterogeneity. PT-based formulations have
been extended to AV--HDV mixed traffic~\cite{adewale2024,ding2021} and to
connected-vehicle compliance~\cite{sharma2019}. A known difficulty with rich
behavioral models is parameter identifiability: many parameters interact and may
not be separately recoverable from short trajectories~\cite{punzo2012}.

\subsection{AV car-following and string stability}
Several studies characterize adaptive-cruise-control and AV CF behavior using
field data~\cite{li2021,makridis2021}, with conflicting conclusions on which
model best fits AV trajectories~\cite{kim2023,bai2024}. Hu et al.~\cite{hu2023}
calibrate IDM on the Waymo dataset and conclude AVs are more conservative (longer
headways, larger jam spacing) and less prone to string instability, while
Alhariqi et al.~\cite{alhariqi2022} introduce an adaptive IDM for mixed autonomy.
String stability---the (non-)amplification of speed perturbations along a
platoon---is a classical property of CF models~\cite{treiber2013,wilson2011}, yet
it is rarely connected back to \emph{empirically calibrated} class-specific
parameters and almost never stratified by traffic regime. This paper makes that
connection explicit.

%================================================================= METHODOLOGY
\section{Methodology}

\subsection{Data}
We use the naturalistic Third-Generation Simulation (TGSIM)
dataset~\cite{ammourah2024}, which captures automated and human-driven vehicles
across several U.S.\ facilities. We use the I-395 (Washington, D.C.) urban
expressway segment and the I-294 (Illinois) segments. A valid car-following
episode requires: duration $>10$~s; bumper-to-bumper gap $<200$~m and $>0$; no
lane change; and a constant preceding-vehicle identity throughout. Followers are
labeled small, large, or autonomous. To ensure comparability, equal numbers of
small and large followers are sampled in the spatial--temporal neighborhood of
each AV episode.

\subsection{Car-following models}
\paragraph{Intelligent Driver Model (IDM).} The acceleration is
\begin{equation}
\dot v_n(t) = a\!\left[1-\left(\frac{v_n(t)}{v_0}\right)^{\delta}
-\left(\frac{s^\ast\!\big(v_n,\dv_n\big)}{s_n(t)}\right)^{2}\right],
\label{eq:idm}
\end{equation}
\begin{equation}
s^\ast(v_n,\dv_n)=s_0+\max\!\left(0,\;v_n T+\frac{v_n\,\dv_n}{2\sqrt{ab}}\right),
\label{eq:sstar}
\end{equation}
with parameters $\Theta_{\mathrm{IDM}}=\{T,a,b,v_0,s_0\}$, where $T$ is the safe
time headway, $a$ the maximum acceleration, $b$ the comfortable deceleration,
$v_0$ the desired speed, and $s_0$ the minimum bumper-to-bumper gap. The
acceleration exponent $\delta$ is fixed at the standard value $\delta=4$; as
shown in Section~\ref{sec:ident}, $\delta$ is not practically identifiable from
short trajectories and ``calibrating'' it injects noise.

\paragraph{Prospect-Theory (PT) model.} The PT model~\cite{hamdar2008rec,
hamdar2015} selects acceleration as a risk-taking decision. The generalized
utility combines a prospect-theoretic acceleration utility and a crash-risk
penalty,
\begin{equation}
U(a;s,v,\dv)=U_{\mathrm{PT}}(a)+U_{\mathrm{crash}}(a;s,v,\dv),
\end{equation}
and the realized acceleration is obtained by a Newton--Raphson solution of the
first-order condition (full equations as in~\cite{hamdar2015}). The calibrated
parameters are
$\Theta_{\mathrm{PT}}=\{T_{\max},\alpha,\beta,W_c,W_m,\gamma\}$: the maximum
anticipation horizon $T_{\max}$, speed-uncertainty coefficient $\alpha$,
intra-driver logit variability $\beta$, accident weight $W_c$, negative-utility
asymmetry $W_m$, and sensitivity exponent $\gamma$.

\paragraph{Optimal-Velocity Relative-Velocity (OVRV) model.} To represent the
optimal-velocity tradition~\cite{bando1995,jiang2001}, we calibrate
\begin{equation}
\dot v_n = \kappa\big(V(s_n)-v_n\big)+\lambda\,(v_{\mathrm{lead}}-v_n),\qquad
V(s)=v_{\max}\,\frac{\tanh\!\big((s-s_c)/s_w\big)+\tanh(s_c/s_w)}{1+\tanh(s_c/s_w)},
\end{equation}
with $\Theta_{\mathrm{OVRV}}=\{\kappa,\lambda,v_{\max},s_c,s_w\}$ (speed
sensitivity $\kappa$, relative-velocity gain $\lambda$, free-flow speed
$v_{\max}$, inflection gap $s_c$, transition width $s_w$); the normalization gives
$V(0)=0$ and $V(\infty)=v_{\max}$.

\paragraph{Gipps safety model.} To represent the safety/collision-avoidance
tradition~\cite{gipps1981}, the next speed is the minimum of a free-acceleration
target $v_{\mathrm{acc}}=v_n+2.5\,a\,\tau(1-v_n/V_0)\sqrt{0.025+v_n/V_0}$ and a
safe-braking target
$v_{\mathrm{safe}}=-b\tau+\sqrt{\max(0,\,b^2\tau^2+b(2(s_n-s_0)-v_n\tau)+v_{\mathrm{lead}}^2)}$,
with $\Theta_{\mathrm{Gipps}}=\{a,b,V_0,s_0,\tau\}$. Because TGSIM is sampled
finely, the $\tau$-ahead Gipps target is applied as an effective acceleration
$(v_{\mathrm{target}}-v_n)/\tau$ and integrated at the data step in closed loop.

\paragraph{ACC-IDM (AV-control surrogate).} As an AV-oriented surrogate we
calibrate the improved IDM with the constant-acceleration heuristic
(CAH)~\cite{kesting2010}, which blends the IDM acceleration $a_{\mathrm{IDM}}$ and
the heuristic $a_{\mathrm{CAH}}$ through a coolness coefficient $c$:
$a=a_{\mathrm{IDM}}$ if $a_{\mathrm{IDM}}\ge a_{\mathrm{CAH}}$, otherwise
$a=(1-c)a_{\mathrm{IDM}}+c\big(a_{\mathrm{CAH}}+b\tanh((a_{\mathrm{IDM}}-a_{\mathrm{CAH}})/b)\big)$,
with the leader acceleration in $a_{\mathrm{CAH}}$ estimated from the observed
leader speed and $\Theta_{\mathrm{ACC}}=\{T,a,b,v_0,s_0,c\}$.

\paragraph{Rationale for four paradigms.} Calibrating models from four distinct
traditions (physics, behavioral, optimal-velocity, and safety, plus an AV
surrogate) tests whether the AV--HDV and regime effects reported below are robust
across modeling assumptions rather than artifacts of a single formulation. All
four use the identical closed-loop simulator and calibration engine described
next.

\subsection{Closed-loop trajectory simulation}
\label{sec:closedloop}
For all models the follower trajectory is integrated forward from the first
observed state $(x_n(0),v_n(0))$ using the observed leader trajectory. Crucially,
the spacing that enters the model is recomputed at each step from the
\emph{simulated} follower position,
\begin{equation}
s_n(t_i)=x_{\mathrm{lead}}(t_i)-x_n(t_i)-\tfrac12 L_{\mathrm{lead}}
-\tfrac12 L_{\mathrm{foll}},
\label{eq:closedloop}
\end{equation}
with a ballistic update $v_{i+1}=\max(0,v_i+a_i\,\Delta t)$,
$x_{i+1}=x_i+v_{i+1}\Delta t$ (PT uses
$x_{i+1}=x_i+v_i\Delta t+\tfrac12 a_i\Delta t^2$). This closed-loop scheme makes
the simulated trajectory---and hence the calibration objective---genuinely
sensitive to the parameters, in contrast to an open-loop scheme that feeds the
observed spacing back into the model and yields optimistic, weakly identifiable
fits. Acceleration is bounded to $[-8,5]~\mathrm{m/s^2}$.

\subsection{Genetic-algorithm calibration}
Parameters are calibrated per leader--follower episode by a genetic algorithm
that minimizes a weighted trajectory error,
\begin{equation}
f(\Theta)=\sum_{j}\Big(\lambda_p\,\big|x_{\mathrm{obs},j}-x_{\mathrm{sim},j}(\Theta)\big|
+\lambda_v\,\big|v_{\mathrm{obs},j}-v_{\mathrm{sim},j}(\Theta)\big|\Big),
\label{eq:fitness}
\end{equation}
with $\lambda_p=1$, $\lambda_v=\tbd{0.5}$. To assess robustness and
identifiability (Section~\ref{sec:ident}) each episode is calibrated from
$N=\tbd{20}$ random restarts.

\subsection{Traffic-regime classification}
\label{sec:regime}
Each episode is assigned a dominant traffic regime from its follower speed
profile $v_n(\cdot)$ using interpretable features: mean speed $\bar v$, speed
standard deviation $\sigma_v$, minimum speed $v_{\min}$, the near-stop fraction
$\phi=\Pr(v<v_{\mathrm{stop}})$, and the speed range $v_{\max}-v_{\min}$. With
$v_{\mathrm{stop}}=3$~m/s, the rules are:
\begin{itemize}
\item \textbf{stop-and-go}: $\phi>0.05$, or ($v_{\min}<v_{\mathrm{stop}}$ and a
large, oscillatory speed range);
\item \textbf{free-flow}: $\bar v\ge 17$~m/s and $\sigma_v\le1.5$~m/s;
\item \textbf{synchronized}: otherwise (congested but flowing).
\end{itemize}
A sliding window labels each instant and contiguous same-regime segments
($\ge10$~s) define the dominant regime. Class comparisons are then performed both
pooled and \emph{within} each regime.

\subsection{Out-of-sample validation and practical identifiability}
\label{sec:ident}
\paragraph{Out-of-sample prediction.} For each episode we calibrate on the first
$70\%$ of the trajectory and predict the remaining $30\%$ in closed loop,
reporting in-sample and out-of-sample position RMSE and their ratio. This guards
against overfitting that pure in-sample error cannot detect.

\paragraph{Practical identifiability.} For a parameter $\theta$ with restart
estimates $\{\hat\theta^{(r)}\}_{r=1}^{N}$, the normalized restart dispersion
\begin{equation}
\mathrm{CV}(\theta)=\frac{\mathrm{std}_r\,\hat\theta^{(r)}}{|\overline{\hat\theta}|}
\label{eq:cv}
\end{equation}
measures how tightly the data pin down $\theta$. We classify
$\mathrm{CV}\le0.10$ as well identifiable, $0.10<\mathrm{CV}\le0.30$ as
moderately identifiable, and $\mathrm{CV}>0.30$ as poorly identifiable. We
complement this with one-at-a-time profiles of the fit around the optimum: a flat
profile indicates a non-identifiable direction.

\subsection{Statistical comparison}
Class differences are tested with Welch's ANOVA (robust to unequal variances) and
Games--Howell post-hoc pairwise tests (robust to unequal variances and sample
sizes), reporting $F$/$t$ statistics, $p$-values, and Hedges' $g$ effect sizes.
The same tests are applied within each regime.

\subsection{Emergent dynamics}
\label{sec:emergent}
\paragraph{Linear string stability.} Linearizing a CF model
$\dot v_n=f(s_n,\dv_n,v_n)$ about a uniform equilibrium $(v_e,s_e)$ (with
$\dv=v_n-v_{n-1}$) gives the speed-to-speed transfer function between consecutive
vehicles
\begin{equation}
H(\omega)=\frac{f_s-i\omega f_{\dv}}{(f_s-\omega^2)-i\omega\,(f_{\dv}+f_v)},
\label{eq:H}
\end{equation}
where $f_s,f_v,f_{\dv}$ are partial derivatives of $f$ at equilibrium. The
platoon is string stable iff $|H(\omega)|\le1$ for all $\omega$; equivalently the
low-frequency criterion
\begin{equation}
\lambda_2=\tfrac12 f_v^2+f_{\dv}f_v-f_s\ \ge\ 0
\label{eq:lambda2}
\end{equation}
must hold~\cite{treiber2013}. For IDM the equilibrium spacing is
$s_e=(s_0+v_eT)/\sqrt{1-(v_e/v_0)^\delta}$ and partials are evaluated
numerically. We report the peak gain $\max_\omega|H(\omega)|$ and the unstable
speed band per class and regime.

\paragraph{Mixed-autonomy platoon simulation.} Using class-specific calibrated
parameters, we simulate a single-lane platoon at AV market-penetration rate (MPR)
$p\in[0,1]$ subject to a stop-and-go lead-vehicle perturbation, averaging over
random class compositions. We report: the empirical string amplification (ratio
of last-to-lead speed-oscillation standard deviation); surrogate safety measures
time-to-collision $\mathrm{TTC}=s/\dv$ (when closing), required deceleration
$\mathrm{DRAC}=\dv^2/(2s)$, time-exposed TTC below $3$~s (TET) and time-integrated
TTC (TIT); and throughput.

%================================================================= RESULTS
\section{Results}

\subsection{Calibration accuracy and out-of-sample validation}
Figure~\ref{fig:fit} shows simulated versus observed position and speed for
representative IDM and PT followers under the closed-loop scheme.
Table~\ref{tab:err} reports mean RMSE and MAE by class for both models, together
with the out-of-sample errors from the $70/30$ temporal split. \tbd{Fill
Table~\ref{tab:err} from the re-run. Note that closed-loop RMSE will be larger
than the previously reported open-loop values; this is the correct, honest
metric. Report the in-/out-of-sample ratio to show the models generalize.}

\begin{figure}[t]
\centering
% Replace these placeholders with your regenerated per-episode comparison plots
% from "Results IDM/comparison_plots/" and "Results PT/comparison_plots/".
\fbox{\begin{minipage}[c][3.2cm][c]{0.46\textwidth}\centering
\tbd{(a) IDM example follower plot}\end{minipage}}\hfill
\fbox{\begin{minipage}[c][3.2cm][c]{0.46\textwidth}\centering
\tbd{(b) PT example follower plot}\end{minipage}}
\caption{Closed-loop simulated (dashed) versus observed (solid) position and
speed for a representative follower: (a) IDM, (b) PT.}
\label{fig:fit}
\end{figure}

\begin{table}[t]
\centering
\caption{Trajectory error by vehicle class (closed-loop), in-sample and
out-of-sample.}
\label{tab:err}
\begin{tabular}{llcccc}
\toprule
Model & Vehicle type & RMSE (m) & MAE (m) & OOS RMSE (m) & OOS/IS ratio\\
\midrule
\multirow{3}{*}{IDM} & Small & \tbd{} & \tbd{} & \tbd{} & \tbd{}\\
 & Large & \tbd{} & \tbd{} & \tbd{} & \tbd{}\\
 & AV    & \tbd{} & \tbd{} & \tbd{} & \tbd{}\\
\midrule
\multirow{3}{*}{PT} & Small & \tbd{} & \tbd{} & \tbd{} & \tbd{}\\
 & Large & \tbd{} & \tbd{} & \tbd{} & \tbd{}\\
 & AV    & \tbd{} & \tbd{} & \tbd{} & \tbd{}\\
\bottomrule
\end{tabular}
\end{table}

\subsection{Practical identifiability}
Table~\ref{tab:ident} reports the median normalized restart dispersion
(Eq.~\ref{eq:cv}) per parameter. \tbd{Fill from
\texttt{analysis/validation\_identifiability.py}.} We expect IDM's $T$, $v_0$,
and (in congested regimes) $s_0$ to be well identified, while several PT
parameters ($T_{\max}$, $\gamma$, $\beta$) are poorly identified. This reframes
the limited PT differentiation observed across classes: where a parameter is not
identifiable from the trajectory, class-level differences are statistically
diluted regardless of any true behavioral difference.

\begin{table}[t]
\centering
\caption{Practical identifiability: median restart-dispersion CV per parameter
(lower $=$ better identified).}
\label{tab:ident}
\begin{tabular}{lccc}
\toprule
Parameter & Small & Large & AV \\
\midrule
IDM $T$  & \tbd{} & \tbd{} & \tbd{}\\
IDM $a$  & \tbd{} & \tbd{} & \tbd{}\\
IDM $b$  & \tbd{} & \tbd{} & \tbd{}\\
IDM $v_0$& \tbd{} & \tbd{} & \tbd{}\\
IDM $s_0$& \tbd{} & \tbd{} & \tbd{}\\
PT $T_{\max}$ & \tbd{} & \tbd{} & \tbd{}\\
PT $\alpha$   & \tbd{} & \tbd{} & \tbd{}\\
PT $\gamma$   & \tbd{} & \tbd{} & \tbd{}\\
\bottomrule
\end{tabular}
\end{table}

\subsection{Pooled and regime-stratified parameter differences}
Table~\ref{tab:params} summarizes calibrated IDM parameters by class (mean $\pm$
SD, median). \tbd{Fill from the re-run \texttt{idm\_parameters\_table.csv}.}
Table~\ref{tab:regdist} reports the regime composition of each class, and
Table~\ref{tab:regstats} the Welch ANOVA of class differences \emph{within} each
regime. \tbd{Fill from \texttt{idm\_regime\_distribution.csv} and
\texttt{idm\_welch\_anova\_by\_regime.csv}.} The central new finding is expected
to be that class differences (e.g., in $T$ and $s_0$) are sharper in the
stop-and-go regime---where car-following is most constrained---than in the pooled
sample, demonstrating that regime pooling masks behavioral distinctions.

\begin{table}[t]
\centering
\caption{Calibrated IDM parameters by vehicle class (closed-loop). Mean $\pm$ SD
(median).}
\label{tab:params}
\begin{tabular}{lccc}
\toprule
Parameter & Small & Large & AV \\
\midrule
$T$ (s)   & \tbd{} & \tbd{} & \tbd{}\\
$a$ (m/s$^2$) & \tbd{} & \tbd{} & \tbd{}\\
$b$ (m/s$^2$) & \tbd{} & \tbd{} & \tbd{}\\
$v_0$ (m/s)   & \tbd{} & \tbd{} & \tbd{}\\
$s_0$ (m)     & \tbd{} & \tbd{} & \tbd{}\\
\bottomrule
\end{tabular}
\end{table}

\begin{table}[t]
\centering
\caption{Regime composition by vehicle class (episode counts).}
\label{tab:regdist}
\begin{tabular}{lccc}
\toprule
Class & Free-flow & Synchronized & Stop-and-go \\
\midrule
Small & \tbd{} & \tbd{} & \tbd{}\\
Large & \tbd{} & \tbd{} & \tbd{}\\
AV    & \tbd{} & \tbd{} & \tbd{}\\
\bottomrule
\end{tabular}
\end{table}

\begin{table}[t]
\centering
\caption{Welch ANOVA of IDM parameters across classes, within each regime
($p$-values; bold $=$ significant at $0.05$).}
\label{tab:regstats}
\begin{tabular}{lccc}
\toprule
Parameter & Free-flow & Synchronized & Stop-and-go \\
\midrule
$T$   & \tbd{} & \tbd{} & \tbd{}\\
$v_0$ & \tbd{} & \tbd{} & \tbd{}\\
$s_0$ & \tbd{} & \tbd{} & \tbd{}\\
\bottomrule
\end{tabular}
\end{table}

\subsection{Emergent string stability}
Using the calibrated class means, Figure~\ref{fig:string} shows the peak
amplification $\max_\omega|H(\omega)|$ and the stability coefficient
$\lambda_2(v_e)$ across equilibrium speeds. With the current (pre-recalibration)
class means, the autonomous class exhibits the \emph{widest} string-unstable speed
band and the highest peak amplification
(\prelim{peak $|H|\approx1.14$, unstable band $\approx1$--$18$~m/s}), compared with
small (\prelim{$1.04$, $1$--$12$~m/s}) and large
(\prelim{$1.00$, $1$--$7$~m/s}) vehicles. \tbd{Regenerate after recalibration via
\texttt{analysis/run\_emergent\_analysis.py --params-csv ...}; also report the
per-regime version using \texttt{--regime}.}

\begin{figure}[t]
\centering
\includegraphics[width=0.9\textwidth]{../Results Emergent/fig_string_stability.png}
\caption{String stability by vehicle class: (left) peak speed-to-speed gain
$\max_\omega|H(\omega)|$ versus equilibrium speed (values $>1$ are string
unstable); (right) low-frequency stability coefficient $\lambda_2$ ($\ge0$
stable).}
\label{fig:string}
\end{figure}

\subsection{Mixed-autonomy safety and throughput}
Figure~\ref{fig:mpr} reports emergent platoon outcomes versus AV market
penetration. With the current class means, string amplification rises with AV
share and crosses the stability threshold near
\prelim{$\sim$80\% penetration}; the mean minimum time-to-collision decreases
(\prelim{$\approx4.8\to3.8$~s}) while throughput increases modestly
(\prelim{$\approx1854\to1969$~veh/h}). The picture is therefore nuanced: the
calibrated AV behavior trades small throughput gains against reduced string
stability and shrinking safety margins. \tbd{Regenerate after recalibration and,
ideally, report this separately for the stop-and-go regime, where the effects are
strongest.}

\begin{figure}[t]
\centering
\includegraphics[width=0.95\textwidth]{../Results Emergent/fig_mpr_sweep.png}
\caption{Emergent outcomes versus AV market penetration: string amplification,
required braking (DRAC), minimum time-to-collision, and throughput.}
\label{fig:mpr}
\end{figure}

%================================================================= DISCUSSION
\section{Discussion}
Three points distinguish this analysis from prior AV--HDV calibration studies.
First, the closed-loop calibration scheme makes the objective genuinely sensitive
to the parameters; combined with the identifiability analysis, it lets us state
\emph{which} parameters are trustworthy rather than reporting all of them at face
value. This directly explains the previously puzzling lack of PT differentiation:
where a parameter is not identifiable, no statistical test can recover a class
difference. Second, regime stratification shows that AV--HDV differences are
condition dependent; treating ``car-following'' as a single regime understates
the differences that matter in congested, safety-critical conditions. Third, the
emergent analysis reframes the policy message: the calibrated AV behavior is not
uniformly beneficial---throughput gains can coincide with reduced string
stability and safety margins, particularly at high penetration and in stop-and-go
conditions. These findings caution against extrapolating ``AVs are conservative
and therefore safer'' from pooled, in-sample parameter comparisons.

%================================================================= CONCLUSION
\section{Conclusion}
We presented a regime-aware, validation-grounded, and consequence-linked
re-analysis of AV versus HDV car-following using TGSIM trajectories. Both IDM and
PT were calibrated with a closed-loop trajectory scheme; parameters were assessed
for out-of-sample accuracy and practical identifiability; class differences were
tested within traffic regimes; and the calibrated parameters were mapped to
emergent string stability and mixed-autonomy safety and throughput. \tbd{Summarize
the final quantitative findings after the re-run: which parameters differ by
class and regime, which are identifiable, the OOS accuracy, and the
penetration/regime at which string stability and safety degrade.} The resulting
class- and regime-specific parameters, together with their identifiability
labels, provide actionable priors for mixed-autonomy simulation and for
lane-management and AV-integration policy. Limitations include the moderate
sample size per class, a single AV fleet with proprietary control logic, and the
linear nature of the string-stability criterion; future work will extend the
analysis to lane-changing and to multi-fleet AV comparisons.

%================================================================= BACK MATTER
\section*{Acknowledgments}
This study was supported by the U.S.\ Department of Transportation (USDOT) Federal
Highway Administration (FHWA) under the project ``New Generation Models for
Connected Autonomous Vehicles (NGM-CAV).'' This work is dedicated to the memory of
Professor Hani Mahmassani.

\section*{Author Contributions}
The authors confirm contribution to the paper as follows: conceptualization,
methodology, validation, analysis, and manuscript preparation as in the original
submission; the v2 methodology (closed-loop calibration, regime stratification,
identifiability, and emergent-dynamics analysis) was led by P.\ Beigi. All authors
reviewed and approved the final manuscript.

\section*{Data Availability}
This study uses the Third-Generation Simulation (TGSIM) data supplied by USDOT.
Analysis code is available in the project repository.

%================================================================= REFERENCES
\begin{thebibliography}{99}
\small
\bibitem{chung2009} E.\ Chung and A.-G.\ Dumont, \emph{Transport Simulation:
Beyond Traditional Approaches}. EPFL Press, 2009.
\bibitem{zheng2023} S.-T.\ Zheng et al., ``A multi-objective calibration
framework for capturing the behavioral patterns of autonomously-driven
vehicles,'' \emph{Transp.\ Res.\ Part C}, vol.\ 152, 104151, 2023.
\bibitem{zheng2016} Z.\ Zheng and M.\ Sarvi, ``Modeling, calibrating, and
validating car following and lane changing behavior,'' \emph{Transp.\ Res.\ Part
C}, vol.\ 71, pp.\ 182--183, 2016.
\bibitem{gunter2020} G.\ Gunter et al., ``Are commercially implemented adaptive
cruise control systems string stable?'' \emph{IEEE Trans.\ Intell.\ Transp.\
Syst.}, vol.\ 22, no.\ 11, pp.\ 6992--7003, 2021.
\bibitem{ammourah2024} R.\ Ammourah et al., ``Introduction to the Third
Generation Simulation Dataset: Data Collection and Trajectory Extraction,''
\emph{Transp.\ Res.\ Rec.}, 2024.
\bibitem{treiber2000} M.\ Treiber, A.\ Hennecke, and D.\ Helbing, ``Congested
traffic states in empirical observations and microscopic simulations,''
\emph{Phys.\ Rev.\ E}, vol.\ 62, no.\ 2, 1805, 2000.
\bibitem{bando1995} M.\ Bando, K.\ Hasebe, A.\ Nakayama, A.\ Shibata, and Y.\
Sugiyama, ``Dynamical model of traffic congestion and numerical simulation,''
\emph{Phys.\ Rev.\ E}, vol.\ 51, no.\ 2, pp.\ 1035--1042, 1995.
\bibitem{jiang2001} R.\ Jiang, Q.\ Wu, and Z.\ Zhu, ``Full velocity difference
model for a car-following theory,'' \emph{Phys.\ Rev.\ E}, vol.\ 64, no.\ 1,
017101, 2001.
\bibitem{gipps1981} P.\ G.\ Gipps, ``A behavioural car-following model for
computer simulation,'' \emph{Transp.\ Res.\ Part B}, vol.\ 15, no.\ 2, pp.\
105--111, 1981.
\bibitem{kesting2010} A.\ Kesting, M.\ Treiber, and D.\ Helbing, ``Enhanced
intelligent driver model to access the impact of driving strategies on traffic
capacity,'' \emph{Phil.\ Trans.\ R.\ Soc.\ A}, vol.\ 368, no.\ 1928, pp.\
4585--4605, 2010.
\bibitem{zhu2018} M.\ Zhu et al., ``Modeling car-following behavior on urban
expressways in Shanghai: A naturalistic driving study,'' \emph{Transp.\ Res.\
Part C}, vol.\ 93, pp.\ 425--445, 2018.
\bibitem{pourabdollah2018} M.\ Pourabdollah et al., ``Calibration and evaluation
of car following models using real-world driving data,'' \emph{IEEE ITSC}, 2018.
\bibitem{punzo2005} V.\ Punzo and F.\ Simonelli, ``Analysis and comparison of
microscopic traffic flow models with real traffic microscopic data,''
\emph{Transp.\ Res.\ Rec.}, vol.\ 1934, pp.\ 53--63, 2005.
\bibitem{punzo2012} V.\ Punzo, M.\ Montanino, and B.\ Ciuffo, ``Do we really need
to calibrate all the parameters? Variance-based sensitivity analysis to simplify
microscopic traffic flow models,'' \emph{IEEE Trans.\ Intell.\ Transp.\ Syst.},
vol.\ 16, no.\ 1, pp.\ 184--193, 2015.
\bibitem{montanino2015} M.\ Montanino and V.\ Punzo, ``Trajectory data
reconstruction and simulation-based validation against macroscopic traffic
patterns,'' \emph{Transp.\ Res.\ Part B}, vol.\ 80, pp.\ 82--106, 2015.
\bibitem{hamdar2008rec} S.\ H.\ Hamdar, M.\ Treiber, H.\ S.\ Mahmassani, and A.\
Kesting, ``Modeling driver behavior as a sequential risk-taking task,''
\emph{Transp.\ Res.\ Rec.}, vol.\ 2088, pp.\ 208--217, 2008.
\bibitem{hamdar2015} S.\ H.\ Hamdar, H.\ S.\ Mahmassani, and M.\ Treiber, ``From
behavioral psychology to acceleration modeling: Calibration, validation, and
exploration of drivers' cognitive and safety parameters in a risk-taking
environment,'' \emph{Transp.\ Res.\ Part B}, vol.\ 78, pp.\ 32--53, 2015.
\bibitem{adewale2024} A.\ Adewale and C.\ Lee, ``Prediction of car-following
behavior of autonomous and human-driven vehicles based on drivers' memory and
cooperation,'' \emph{Transp.\ Res.\ Rec.}, vol.\ 2678, no.\ 6, pp.\ 248--266, 2024.
\bibitem{ding2021} S.\ Ding et al., ``An extended car-following model in
connected and autonomous vehicle environment,'' \emph{J.\ Adv.\ Transp.}, 2739129,
2021.
\bibitem{sharma2019} A.\ Sharma et al., ``Modelling car-following behaviour of
connected vehicles with a focus on driver compliance,'' \emph{Transp.\ Res.\ Part
B}, vol.\ 126, pp.\ 256--279, 2019.
\bibitem{li2021} T.\ Li et al., ``Car-following behavior characteristics of
adaptive cruise control vehicles based on empirical experiments,'' \emph{Transp.\
Res.\ Part B}, vol.\ 147, pp.\ 67--91, 2021.
\bibitem{makridis2021} M.\ Makridis et al., ``OpenACC: An open database of
car-following experiments to study the properties of commercial ACC systems,''
\emph{Transp.\ Res.\ Part C}, vol.\ 125, 103047, 2021.
\bibitem{kim2023} B.\ Kim and K.\ P.\ Heaslip, ``Identifying suitable
car-following models to simulate automated vehicles on highways,'' \emph{Int.\ J.\
Transp.\ Sci.\ Technol.}, vol.\ 12, no.\ 2, pp.\ 652--664, 2023.
\bibitem{bai2024} J.\ Bai, S.\ Mao, and J.\ J.\ Lee, ``A novel car-following
model for adaptive cruise control vehicles using enhanced IDM,'' \emph{Transp.\
Lett.}, 2024.
\bibitem{hu2023} X.\ Hu, Z.\ Zheng, D.\ Chen, and J.\ Sun, ``Autonomous vehicle's
impact on traffic: Empirical evidence from Waymo Open Dataset and implications
from modelling,'' \emph{IEEE Trans.\ Intell.\ Transp.\ Syst.}, vol.\ 24, no.\ 6,
pp.\ 6711--6724, 2023.
\bibitem{alhariqi2022} A.\ Alhariqi, Z.\ Gu, and M.\ Saberi, ``Calibration of the
IDM with adaptive parameters for mixed autonomy traffic using experimental
trajectory data,'' \emph{Transportmetrica B}, vol.\ 10, no.\ 1, pp.\ 421--440, 2022.
\bibitem{treiber2013} M.\ Treiber and A.\ Kesting, \emph{Traffic Flow Dynamics:
Data, Models and Simulation}. Springer, 2013.
\bibitem{wilson2011} R.\ E.\ Wilson and J.\ A.\ Ward, ``Car-following models:
fifty years of linear stability analysis---a mathematical perspective,''
\emph{Transp.\ Plan.\ Technol.}, vol.\ 34, no.\ 1, pp.\ 3--18, 2011.
\end{thebibliography}

\end{document}
